\appendix

\section{Reader Prompts}
\label{app:prompts}
The reader receives a system message and a user message through its chat template.
Decoding is greedy, with at most 32 new tokens and a stop at the first newline.
Prompt texts are copied from the frozen code; the hash is the first 12 hexadecimal characters of the SHA-256 of the system text, as logged with every output.

\noindent\textbf{A1: main prompt} (hash \texttt{c50efcc15cc6}).\par
\begin{lstlisting}
[system]
You are a question answering assistant. Answer the question using the passages provided. If the passages do not contain the answer, give your best short answer anyway. Reply with the answer only, as a short phrase of a few words, with no explanation.
[user]
Passages:
[1] {title 1}
{text 1}

[2] {title 2}
{text 2}
...

Question: {question}
Answer:
\end{lstlisting}

\noindent\textbf{A2: closed-book prompt} (hash \texttt{c22f911b0ee0}).\par
\begin{lstlisting}
[system]
You are a question answering assistant. Reply with the answer only, as a short phrase of a few words, with no explanation.
[user]
Question: {question}
Answer:
\end{lstlisting}

\noindent\textbf{A3: abstention prompt} (hash \texttt{cdc86e94862c}).\par
The user message is identical to A1.
\begin{lstlisting}
[system]
You are a question answering assistant. Answer the question using only the passages provided. If the passages do not contain the answer, reply exactly: I don't know. Otherwise reply with the answer only, as a short phrase of a few words, with no explanation.
\end{lstlisting}
An output counts as ``I don't know'' if, ignoring case and leading spaces, it begins with ``i don't know'', ``i dont know'' or ``i do not know'' (straight or curly apostrophe); such outputs score EM 0.

\section{Judge Prompt}
\label{app:judge}
The judge receives one user message and no system message.
Decoding is greedy with at most 8 new tokens.
We parse the output with \verb|^\s*(SUPPORTED|\verb/|NOT_SUPPORTED)\b/ and do not retry.
A pre-registered rule switched to a few-shot variant only if more than 2\% of zero-shot outputs failed to parse; 0 of 3,999 failed, so all verdicts use the zero-shot prompt (hash \texttt{da33361c8520}).
Outputs that are not answers (for example an empty output at the token limit) are not judged.
\begin{lstlisting}
You will be given PASSAGES, a QUESTION, and an ANSWER.
Decide whether the ANSWER is fully supported by the PASSAGES.
An answer is SUPPORTED only if the passages state it or directly imply it.
Do not use outside knowledge. Do not judge whether the answer is true in the real world.

PASSAGES:
[1] {title 1}
{text 1}
...

QUESTION: {question}
ANSWER: {answer}

Reply with exactly one word: SUPPORTED or NOT_SUPPORTED.
\end{lstlisting}

\section{Additional Tables}
\label{app:tables}
Table~\ref{tab:TC1} lists every pairwise comparison behind the significance statements in Section~\ref{sec:results}.
Table~\ref{tab:TC2} gives the top-$k$ ablation, Table~\ref{tab:TC3} the abstention prompt and Table~\ref{tab:TC4} the pinned model revisions.
% GENERATED by src/paper/make_latex_tables.py -- do not edit by hand
\begin{table}[t]
\centering
\scriptsize
\setlength{\tabcolsep}{4pt}
\begin{adjustbox}{max width=\linewidth}
\begin{tabular}{lrrrr}
\toprule
A -- B & H@5 & nDCG@10 & EM & SUPP. \\
\midrule
BM25 -- MiniLM-L6-v2 & $-$30.1$^*$ & $-$24.3$^*$ & $-$11.2$^*$ & $-$9.6$^*$ \\
BM25 -- bge-small & $-$36.6$^*$ & $-$30.3$^*$ & $-$12.4$^*$ & $-$10.8$^*$ \\
BM25 -- bge-base & $-$40.7$^*$ & $-$34.9$^*$ & $-$13.3$^*$ & $-$11.2$^*$ \\
BM25 -- bge-large & $-$43.3$^*$ & $-$36.3$^*$ & $-$14.4$^*$ & $-$12.2$^*$ \\
BM25 -- e5-base-v2 & $-$45.3$^*$ & $-$39.0$^*$ & $-$15.0$^*$ & $-$12.6$^*$ \\
BM25 -- Qwen3-Emb-0.6B & $-$38.1$^*$ & $-$32.3$^*$ & $-$13.2$^*$ & $-$10.6$^*$ \\
MiniLM-L6-v2 -- bge-small & $-$6.5$^*$ & $-$6.0$^*$ & $-$1.2 & $-$1.2 \\
MiniLM-L6-v2 -- bge-base & $-$10.6$^*$ & $-$10.6$^*$ & $-$2.2 & $-$1.6 \\
MiniLM-L6-v2 -- bge-large & $-$13.2$^*$ & $-$12.0$^*$ & $-$3.2$^*$ & $-$2.6 \\
MiniLM-L6-v2 -- e5-base-v2 & $-$15.3$^*$ & $-$14.7$^*$ & $-$3.8$^*$ & $-$3.0 \\
MiniLM-L6-v2 -- Qwen3-Emb-0.6B & $-$8.0$^*$ & $-$8.0$^*$ & $-$2.0 & $-$1.0 \\
bge-small -- bge-base & $-$4.1$^*$ & $-$4.6$^*$ & $-$0.9 & $-$0.4 \\
bge-small -- bge-large & $-$6.7$^*$ & $-$6.0$^*$ & $-$2.0 & $-$1.4 \\
bge-small -- e5-base-v2 & $-$8.7$^*$ & $-$8.7$^*$ & $-$2.6$^*$ & $-$1.8 \\
bge-small -- Qwen3-Emb-0.6B & $-$1.5 & $-$2.0$^*$ & $-$0.8 & +0.2 \\
bge-base -- bge-large & $-$2.6$^*$ & $-$1.4$^*$ & $-$1.0 & $-$1.0 \\
bge-base -- e5-base-v2 & $-$4.6$^*$ & $-$4.1$^*$ & $-$1.6 & $-$1.4 \\
bge-base -- Qwen3-Emb-0.6B & +2.6$^*$ & +2.6$^*$ & +0.1 & +0.6 \\
bge-large -- e5-base-v2 & $-$2.0$^*$ & $-$2.7$^*$ & $-$0.6 & $-$0.4 \\
bge-large -- Qwen3-Emb-0.6B & +5.2$^*$ & +4.0$^*$ & +1.2 & +1.6 \\
e5-base-v2 -- Qwen3-Emb-0.6B & +7.2$^*$ & +6.7$^*$ & +1.8 & +2.0 \\
\bottomrule
\end{tabular}
\end{adjustbox}
\caption{All pairwise differences A$-$B in percentage points (joined queries; SUPPORTED on the judge subsample, unvalidated). $^*$: significant after Holm over 21 pairs (exact McNemar for binary measures, Wilcoxon for nDCG@10).}
\label{tab:TC1}
\end{table}

% GENERATED by src/paper/make_latex_tables.py -- do not edit by hand
\begin{table}[t]
\centering
\scriptsize
\setlength{\tabcolsep}{4pt}
\begin{adjustbox}{max width=\linewidth}
\begin{tabular}{lrrrr}
\toprule
Method & $k{=}1$ & $k{=}3$ & $k{=}5$ & $k{=}10$ \\
\midrule
BM25 & 23.3 {\scriptsize[21.6, 25.1]} & 31.8 {\scriptsize[29.9, 33.7]} & 36.0 {\scriptsize[34.0, 38.0]} & 40.0 {\scriptsize[37.9, 42.0]} \\
MiniLM-L6-v2 & 35.2 {\scriptsize[33.2, 37.2]} & 44.2 {\scriptsize[42.1, 46.2]} & 47.1 {\scriptsize[45.1, 49.2]} & 50.2 {\scriptsize[48.1, 52.3]} \\
bge-large & 42.1 {\scriptsize[40.0, 44.2]} & 49.3 {\scriptsize[47.2, 51.3]} & 50.3 {\scriptsize[48.2, 52.4]} & 52.0 {\scriptsize[49.9, 54.0]} \\
\midrule
Mean prompt tokens (R4) & 210 & 472 & 738 & 1,414 \\
\bottomrule
\end{tabular}
\end{adjustbox}
\caption{Top-$k$ ablation: EM (\%) with 95\% CI, n = 2,255. Adjacent steps are tested with exact McNemar, Holm over the 9 steps.}
\label{tab:TC2}
\end{table}

% GENERATED by src/paper/make_latex_tables.py -- do not edit by hand
\begin{table}[t]
\centering
\scriptsize
\setlength{\tabcolsep}{4pt}
\begin{adjustbox}{max width=\linewidth}
\begin{tabular}{lrrrrrr}
\toprule
Cond. & IDK & IDK$\mid$hit & IDK$\mid$miss & EM main & EM IDK & $p$ \\
\midrule
BM25 & 21.5 & 1.3 & 34.7 & 39.0 & 38.0 & 0.688 \\
MiniLM-L6-v2 & 11.5 & 1.6 & 27.6 & 47.0 & 45.5 & 0.453 \\
bge-large & 4.5 & 2.4 & 13.9 & 49.5 & 46.5 & 0.07 \\
Gold & 6.0 & -- & -- & -- & -- & -- \\
\bottomrule
\end{tabular}
\end{adjustbox}
\caption{IDK prompt on the first 200 queries (\%). IDK: share of ``I don't know'' answers. $p$: exact McNemar, EM under the IDK prompt vs the main prompt (uncorrected).}
\label{tab:TC3}
\end{table}

% GENERATED by src/paper/make_latex_tables.py -- do not edit by hand
\begin{table}[t]
\centering
\scriptsize
\setlength{\tabcolsep}{4pt}
\begin{adjustbox}{max width=\linewidth}
\begin{tabular}{lll}
\toprule
Role & Model & Revision \\
\midrule
MiniLM-L6-v2 & \texttt{sentence-transformers/all-MiniLM-L6-v2} & \texttt{1110a243fdf4} \\
bge-small & \texttt{BAAI/bge-small-en-v1.5} & \texttt{5c38ec7c405e} \\
bge-base & \texttt{BAAI/bge-base-en-v1.5} & \texttt{a5beb1e3e68b} \\
bge-large & \texttt{BAAI/bge-large-en-v1.5} & \texttt{d4aa6901d3a4} \\
e5-base-v2 & \texttt{intfloat/e5-base-v2} & \texttt{f52bf8ec8c71} \\
Qwen3-Emb-0.6B & \texttt{Qwen/Qwen3-Embedding-0.6B} & \texttt{97b0c614be4d} \\
Reader & \texttt{Qwen/Qwen2.5-7B-Instruct-AWQ} & \texttt{b25037543e93} \\
Judge & \texttt{mistralai/Mistral-7B-Instruct-v0.3} & \texttt{c170c708c41d} \\
\bottomrule
\end{tabular}
\end{adjustbox}
\caption{Pinned model revisions (Hugging Face commit, first 12 characters).}
\label{tab:TC4}
\end{table}


\section{Pre-registration Record}
\label{app:prereg}
The design document and its addenda are dated and versioned in our repository.
Each addendum was written before the stage it governs.
\begin{itemize}
\item \textbf{v0.2 (26 Sep 2026).} Research questions, retrieval methods, reader, judge, metrics, statistical tests, the transfer-ratio reporting rule and the pilot plan.
\item \textbf{v0.2.1 (26 Sep 2026).} Follow-up rules: the abstention check (F1), a few-shot fallback for the judge (F2), a CPU fallback for Qwen3-Emb-0.6B (F3) and stopping rules for the GPU budget (F4.1--F4.4).
\item \textbf{v0.2.2 (26 Sep 2026).} Deviation D1: the join gate changed from a join rate of at least 0.70 to at least 1,000 joined queries, recorded after the pilot join and before any retrieval or answer outcome; the join diagnostics D1a and D1b.
\item \textbf{v0.2.3 (27 Sep 2026).} Verification of pilot part A; main-run adjustments M1--M5, including the sub-corpus membership rule (M1); deviation D2, dropping the RAGAS cross-check.
\item \textbf{v0.2.4 (27--29 Sep 2026).} Verification of pilot part B, the main-run plan, execution records for S1, S3a, S3b and S4, the clarification of F4.2 before S3a, the failed first S3a attempt, the S1 cost-estimate error and the decision to cut the dimension ablation (A3).
\end{itemize}
Pilot numbers were used only for planning (power and budget) and are not reported as results.
